
Reinforcement Learning from Human Feedback (RLHF) trains language models against a reward model proxy rather than direct human judgment. As optimization pressure increases, this proxy diverges from held-out human preference --- a phenomenon known qualitatively as reward hacking. We ask: how fast does this divergence grow, and does the growth rate replicate across independent experimental settings?

We define the \textit{calibration-alignment divergence curve} as the regression slope of the normalized proxy-gold gap (RM score minus held-out human preference, both on [0,1] scale) on KL divergence budget. Analyzing published RLHF experimental data from two independent sources, we find this slope is $\beta$ = 0.1433 $nat^{-1}$ in Coste et al. [2023] ($R^2$ = 0.958, p = 8.89 $\times 10^{-7})$ and $\beta$ = 0.1599 $nat^{-1}$ in Gao et al. [2023] (p = 0.0025) --- a cross-dataset slope ratio of 1.116, indicating near-identical divergence rates despite different model families and scales.

These results provide the first regression characterization of the \textit{normalized divergence gap} (RM\_norm $-$ gold\_preference) as a function of KL budget with cross-dataset slope comparison, framing reward hacking as an empirical instance of bidirectional alignment tension. We propose the normalized divergence gap as a standard cross-study comparison instrument computable from existing RLHF evaluation infrastructure.

